\documentclass[10pt,a4paper]{amsart}
\usepackage[margin=19mm]{geometry}
\usepackage{amsmath,amssymb,mathtools}
\usepackage[scr=rsfs,cal=euler]{mathalfa}
\usepackage{xfrac}
\usepackage[unicode,hidelinks,bookmarks=false]{hyperref}
\newcommand{\ie}{i.e.\ }
\newcommand{\cf}{cf.\ }
\newcommand{\resp}{respectively\ }
\newcommand{\Subsection}{\subsection}
\providecommand{\nobreakdash}{\nobreak\hbox{-}}
\setlength{\emergencystretch}{2em}
\multlinegap0pt
\usepackage{amssymb}
\usepackage{mathtools}
\usepackage{xcolor}
\usepackage[normalem]{ulem}
\usepackage{tikz}
\newtheorem{lemma}{Lemma}
\newcommand\RR{\ensuremath{\mathbb{R}}}
\title{Existence, uniqueness, stability, and monotonicity of traveling waves for repulsion/attraction chemotaxis models with logistic type source}
\author{Wenxian Shen}
\date{Source: arXiv:2605.04401v1 (CC BY 4.0)}
\begin{document}
\maketitle

\noindent\emph{Source and licence.} Existence, uniqueness, stability, and monotonicity of traveling waves for repulsion/attraction chemotaxis models with logistic type source. Version arXiv:2605.04401v1 (CC BY 4.0), distributed by arXiv under the Creative Commons Attribution 4.0 International licence, http://creativecommons.org/licenses/by/4.0/, as recorded in the arXiv licence metadata for this exact version and shown on the abstract page https://arxiv.org/abs/2605.04401v1. Full licence notice: \texttt{licenses/arxiv-2605.04401-CC-BY-4.0.txt}.\par\smallskip

\noindent\textit{Editorial scope.} Counted unit: the article's lemma estimate-lm2 (source0.tex lines 2588--2598). The article's complete original proof of it is delivered with it (source0.tex lines 2600--2696, one \texttt{proof} environment of 97 lines). No mathematics is written, repaired or re-derived: the statement and the proof are the article's own lines, in the article's own notation, and no retained line is substituted or re-flowed.  Retained uncounted as the source-local context the proof needs: lines 491--494; lines 520--523; lines 1538--1544; lines 1602--1627; lines 1860--1878; lines 2325--2371.  Nothing was omitted from any retained span, so the package carries no omission record.  No retained line contains a citation, so the wrapper carries no bibliography and no bibliography entry.  No retained line contains a figure environment, an image-inclusion command or drawing code, so no image is delivered and the package declares no asset dependency.  Editorial formatting only, in addition: an AMS \texttt{amsart} preamble that restates the article's own declarations and macros used by the retained lines, namely the environments \texttt{lemma} on separate counters, together with the standard packages those lines need.  The retained environments print in this document as Lemma 1 in order of appearance, whereas the article prints the same items with the numbering of its own full document.  The complete native source of the article as distributed in the arXiv e-print for this exact version is bundled unchanged as \texttt{sources/199858/source0.tex}.
\begin{align}
\label{cond-on-c-eq1}
c>c^*_{\chi,m,\gamma}:=\max\Big\{ \frac{1}{m}+m, \frac{1}{\sqrt{m\gamma |\chi|+\gamma^2 |\chi|+\gamma^2}} + \sqrt {m\gamma |\chi| +\gamma^2 |\chi|+\gamma^2}\Big\},
\end{align}

\begin{equation}
\label{chi-star-eq}
\chi^*=\chi^*(m,\alpha,\gamma):=\min\Big\{1, \frac{2m+2\gamma}{m^2+m+2\gamma}\Big\}.
\end{equation}

\begin{equation}
\label{main-moving-eq1}
\begin{cases}
u_t=u_{xx}+cu_x-\chi(u^m v_x)_x +u(1-u^\alpha),\quad &x\in\RR\cr
0=v_{xx}-v+u^\gamma,\quad&x\in\RR.
\end{cases}
\end{equation}

\begin{lemma}[Super-solutions]
\label{prelim-super-lm} 
Assume that one of the following sets of conditions is satisfied:
\begin{itemize}
\item[(i)]  $\alpha\le m+\gamma -1$,  $\chi\le 0$, 
   $0<\kappa <1$ is such that $c:=\kappa+\frac{1}{\kappa}>c_{\chi,m,\gamma}^*$,
 and
$1\le M\le M_{\chi,\kappa,m,\gamma}:=\frac{1}{\kappa \sqrt {\gamma^2+\gamma^2|\chi|+m\gamma |\chi|}}(>1)$, where $c^*_{\chi,m,\gamma}$ is as in  \eqref{cond-on-c-eq1}.
  
\item[(ii)] $\alpha=m+\gamma-1$,    $0\le \chi<\chi^*(m,\alpha,\gamma)$,  
   $0<\kappa <1$,  and   $M\ge M_\chi:=\big( \frac{1}{1-\chi}\big)^{1/\alpha}$, where $\chi^*(m,\alpha,\gamma)$ is as in \eqref{chi-star-eq}.
\end{itemize}
Then 
for any $u\in \mathcal{E}_{\kappa,M,T}$, the following hold:
\begin{itemize}
\item[(1)]  $W= e^{-\kappa x}$ is a super-solution of
$$
W_t=\mathcal{A}(W;u), \quad t\in (0,T),\, \,\, x\ge -\frac{\ln M}{\kappa}.
$$

\item[(2)]   $W=M$ is a super-solution of
$$
W_t=\mathcal{A}(W;u), \quad t\in (0,T),\,\,\, x\in\RR.
$$
\end{itemize}
\end{lemma}

\begin{lemma}[Sub-solutions]
\label{prelim-sub-lm}
Assume that  $0<\kappa<1$,   $\tilde \kappa$ is such that $0<\kappa<\tilde \kappa \le \min\{ (1+\alpha)\kappa,  m\kappa +1/2,  1\}$, and $M\ge 1$.
Then 
for any $u\in \mathcal{E}_{\kappa,M,T}$, the following hold:
\begin{itemize}
\item[(1)] For $D> D_{M,\kappa,\tilde\kappa,\chi,m,\gamma}$, 
{$W=U_{\kappa,\tilde\kappa, D}(x)$}  is a sub-solution of
$$
W_t= \mathcal{A}(W;u)\quad  t\in (0,T),\quad t\in (0,T),\,\,  x>x^-_{\kappa,\tilde\kappa,D}.
$$


\item[(2)]   For $0<d\le  d_{\kappa,\tilde \kappa,D,\chi}:=\min\Big\{\frac{1}{1+|\chi|}, \Big(\frac{\kappa}{\tilde\kappa D}\Big)^{\frac{\kappa}{\tilde \kappa -\kappa}}\Big(1-\frac{\kappa}{\tilde \kappa}\Big)\Big\}$, $W=d$ is a sub-solution of
$$
W_t= \mathcal{A}(W;u),\quad t\in (0,T),\,\,  x\in\RR.
$$ 
\end{itemize}
\end{lemma}

\begin{lemma}
\label{estimate-lm1}
Suppose that $c>2$, and  $(u(x),v(x))$ is  s stationary  solution of \eqref{main-moving-eq1}  connecting $(1,1)$ and $(0,0)$  and $0<u(x)\le \min\{M_\chi, e^{-\kappa x}\}$, where $k=\frac{c-\sqrt{c^2-4}}{2}$.

\begin{itemize}
\item[(1)] There hold
\begin{equation}
\label{estimate-eq0}
|v(x)|\le  M_\chi^\gamma 
 \,\,\, {\rm and}\,\,\,
|v_x(x)|\le  M_\chi^\gamma  \quad \forall\, x\in\RR.
\end{equation}
\end{itemize}

If, in addition,   {$c>\gamma+\frac{1}{\gamma}$,}  which is equivalent to $\kappa<\frac{1}{\gamma}$, then
\begin{equation}
\label{estimate-eq1}
|v(x)|\le  \min\big\{M_\chi^\gamma , \frac{1}{1-\kappa^2 \gamma^2}e^{-\kappa\gamma  x}\big \}
 \,\,\, {\rm and}\,\,\,
|v_x(x)|\le  \min\{M_{\chi}^\gamma ,  \frac{1}{1-\kappa^2 \gamma^2}  e^{-\kappa\gamma  x}\}  \quad \forall\, x\in\RR.
\end{equation}

\begin{itemize}
\item[(2)] $\lim_{x\to\pm \infty} u_x(x)=0$.

\item[(3)] If {$c> m|\chi| M^{m+\gamma-1}_\chi$,}  then  
\begin{equation}
\label{estimate-eq2}
 -\frac{|\chi| M_{\chi}^{m+\gamma}+M_{\chi}}{c-m|\chi| M^{m+\gamma-1} _{\chi} } \,   \le\,  u_x(x)\,  \le  \,  \frac{|\chi| M_\chi^{m+\gamma} +M_\chi(M_\chi^\alpha -1)}{c- m |\chi| M^{m+\gamma-1}_\chi }  \quad \forall\, x\in\RR.
\end{equation}


\item[(4)]  If {$c>\max\big\{\gamma+\frac{1}{\gamma}, m|\chi| M^{m+\gamma-1}_\chi\big\}$,}  then 
\begin{align}
\label{estimate-eq3}
|u_x(x)|\le  \Big(1+\frac{2}{c}\Big) \Big(M_{1,\kappa,\chi,m,\alpha,\gamma}  e^{-\kappa x}+M_{2,\kappa,\chi,m,\alpha,\gamma}e^{-\gamma\kappa x}\Big)\quad \forall\, x\in\RR,
\end{align}
where
$$
M_{1,\kappa,\chi,m,\alpha,\gamma}:= 1+2 |\chi| M^{m+\gamma-1}_\chi+M_\chi^{\alpha}.
$$
and
$$
M_{2, \kappa,\chi,m,\alpha,\gamma}:=\frac{|\chi| m M_{\chi}^{m-1} (|\chi| M_{\chi}^{m+\gamma}+M_\chi^{\alpha+1}) }{(c- m|\chi| M^{m+\gamma-1}_{\chi})(1-\gamma^2\kappa^2)}.
$$
\end{itemize}
\end{lemma}

\begin{lemma}
\label{estimate-lm2}Assume that $c>\max\big\{\gamma+\frac{1}{\gamma}, m|\chi| M^{m+\gamma-1}_\chi\big\}$.   Suppose that  $(u(x),v(x))$ is  s stationary  solution of \eqref{main-moving-eq1}  connecting $(1,1)$ and $(0,0)$  and $0<u(x)\le \min\{M_\chi, e^{-\kappa x}\}$, where $\kappa=\frac{c-\sqrt{c^2-4}}{2}$.
Then
$$
\big|\frac{u_x(x)}{u(x)}\big|\le  \widetilde M_{c,\chi,m,\alpha,\gamma},
$$
where
$$
\widetilde M_{c,\chi,m,\alpha,\gamma}:= \frac{1}{2}\left( c+|\chi| m M_{\chi}^{m+\gamma-1}+\sqrt{ (c+|\chi| m M_{\chi}^{m+\gamma-1})^2 +4|\chi| M_{\chi}^{m+\gamma-1}+4M_{\chi}^\alpha}\right).
$$
\end{lemma}

\begin{proof}
First, let $w(x)=\frac{u_x(x)}{u(x)}$.
Note that
$$
u_{xx}+c u_x(x)-\chi m u^{m-1} u_x(x) v_x(x)-\chi u^m(x) (v(x)-u^\gamma (x))+u(x) (1-u^\alpha(x))=0.
$$
Hence
\begin{equation}
\label{proof-estimate-3-eq1}
\frac{u_{xx}}{u}+c w -\chi m u^{m-1} v_x w -\chi u^{m-1} (v-u^\gamma)+(1-u^\alpha(x))=0.
\end{equation}
By Lemmas \ref{prelim-super-lm} and \ref{prelim-sub-lm}, and   Lemma \ref{estimate-lm1}(2) and (4), 
$$
\lim_{x\to -\infty} w(x)=0,\quad \limsup_{x\to\infty}|w(x)|<\infty.
$$
Hence 
$$
\|w\|_\infty<\infty.
$$
Then by \eqref{proof-estimate-3-eq1},
$$
\|\frac{u_{xx}}{u}\|_\infty<\infty.
$$

Next, if there is $x_0\in\RR$ such that
$$
|w(x_0)|=\|w\|_\infty,
$$
then
$$
w_x(x_0)=\frac{u_{xx}(x_0)}{u(x_0)}-w^2 (x_0)=0.
$$
This together with \eqref{proof-estimate-3-eq1} implies that
$$
w(x_0)^2+\Big(c-\chi m u^{m-1}(x_0)v_x(x_0)\big) w(x_0) -\chi u^{m-1}(x_0) (v(x_0)-u^\gamma(x_0))+(1-u^\alpha(x)(x_0))=0
$$
Let 
$$a(x_0)=c-\chi m u^{m-1}(x_0)v_x(x_0),
$$
and
$$
b(x_0)=-\chi u^{m-1}(x_0) (v(x_0)-u^\gamma(x_0))+(1-u^\alpha(x)(x_0)).
$$
Then 
$$
w(x_0)= \frac{-a(x_0)\pm \sqrt{(a(x_0))^2- 4 b(x_0)}}{2}.
$$
Hence
\begin{align*}
|w(x_0)|&\le\frac{|a(x_0)| +\sqrt{|a(x_0)|^2+4|b(x_0)|}}{2}\\
&\le \frac{1}{2}\left( c+|\chi| m M_{\chi}^{m+\gamma-1}+\sqrt{ (c+|\chi| m M_{\chi}^{m+\gamma-1})^2 +4|\chi| M_{\chi}^{m+\gamma-1}+4M_{\chi}^\alpha}\right)\\
&=\widetilde M_{c,\chi,m,\alpha,\gamma}.
\end{align*}



If $w(x_0)\not = \|w\|_\infty$ for any $x_0\in\RR$, then there is $x_n\to\infty$ such that
$$
\lim_{n\to\infty} |w(x_n)|=\|w\|_\infty.
$$
Note that
$$
w_x=\frac{u_{xx}}{u}-w^2.
$$
Hence $\|w_x\|_\infty<\infty$. This together with \eqref{proof-estimate-3-eq1} implies that
$\|\big(\frac{u_{xx}}{u}\big)_x\|_\infty<\infty$ and then
$$
\|w_{xx}\|_\infty=\|\big(\frac{u_{xx}}{u}\big)_x-2w w_x\|_\infty<\infty.
$$


  Without loss of generality, we may assume that
$$
w^*(x):=\lim_{x\to\infty} w(x+x_n) \quad {\rm and} \quad w^*_x(x)=\lim_{n\to\infty} w_x(x+x_n)
$$
exist locally uniformly in $x\in\RR$.  Clearly,
$$
|w^*(0)|=\|w^*\|_\infty=\|w\|_\infty,
$$
and 
$$
\lim_{n\to\infty} \frac{u_{xx}(x_n)}{u(x_n)}=\lim_{n\to\infty} \big(w^2(x_n)+w_x(x_n)\big)=(w^*(0))^2+w_x^*(0)=(w^*(0))^2.
$$
Then  by \eqref{proof-estimate-3-eq1}, it holds that
$$
(w^*(0))^2+cw^*(0) +1=0.
$$
This implies that
$$
\|w\|_\infty=|w^*(0)|\le \frac{c+\sqrt{c^2-4}}{2}.
$$
Note that
$$
\frac{c+\sqrt{c^2-4}}{2}\le \widetilde M_{c,\chi,m,\alpha,\gamma}.
$$
The lemma is thus proved.
\end{proof}

\end{document}
